\textbf{Phần 1. Trắc nghiệm lựa chọn}
\setcounterpageref{ex}{0}

\begin{ex} %[2D2V14]
Tính giá trị của biểu thức: $18 - 9 + 8 - 7$.
\choice
{$9$}
{\True $10$}
{$11$}
{$8$}
\loigiai{
$18 - 9 = 9$; $9 + 8 = 17$; $17 - 7 = 10$.
\textbf{Đáp án của câu hỏi là B}
}
\end{ex}

\begin{ex} %[2D2V14]
Một đàn cừu có $16$ con. Số cừu trắng nhiều hơn số cừu đen $4$ con. Hỏi có bao nhiêu con cừu trắng?
\choice
{$8$ con}
{$9$ con}
{$12$ con}
{\True $10$ con}
\loigiai{
Số cừu trắng là: $(16 + 4) : 2 = 10$ (con).
\textbf{Đáp án của câu hỏi là D}
}
\end{ex}

\begin{ex} %[2D2V14]
Tìm $x$ biết: $15 - x = 14 - 8$.
\choice
{\True $9$}
{$8$}
{$10$}
{$7$}
\loigiai{
Vế phải: $14 - 8 = 6$. Ta có: $15 - x = 6 \Rightarrow x = 15 - 6 = 9$.
\textbf{Đáp án của câu hỏi là A}
}
\end{ex}

\begin{ex} %[2D2V14]
Mẹ đi chợ mua $1\text{ chục}$ quả trứng gà và $7$ quả trứng vịt. Mẹ đã dùng hết $8$ quả trứng. Hỏi mẹ còn lại bao nhiêu quả trứng?
\choice
{$8$ quả trứng}
{$10$ quả trứng}
{\True $9$ quả trứng}
{$7$ quả trứng}
\loigiai{
$1\text{ chục} = 10$. Tổng số trứng: $10 + 7 = 17$ (quả).\nCòn lại: $17 - 8 = 9$ (quả).
\textbf{Đáp án của câu hỏi là C}
}
\end{ex}

\begin{ex} %[2D2V14]
Điền số thích hợp: $8 + 7 - \dots = 6$.
\choice
{$8$}
{\True $9$}
{$10$}
{$7$}
\loigiai{
$8 + 7 = 15$. Ta có: $15 - 6 = 9$. Vậy điền số 9.
\textbf{Đáp án của câu hỏi là B}
}
\end{ex}

\begin{ex} %[2D2V14]
Ba bạn An, Bình, Cư thi ném bóng. An ném trúng $7$ quả, Bình ném trúng $8$ quả, Cư ném trúng ít hơn tổng số bóng của An và Bình là $6$ quả. Cư ném trúng bao nhiêu quả bóng?
\choice
{$8$ quả bóng}
{$10$ quả bóng}
{$7$ quả bóng}
{\True $9$ quả bóng}
\loigiai{
Tổng bóng của An và Bình: $7 + 8 = 15$ (quả).\nCư ném trúng: $15 - 6 = 9$ (quả bóng).
\textbf{Đáp án của câu hỏi là D}
}
\end{ex}

\begin{ex} %[2D2V14]
Tìm số tự nhiên lớn nhất có một chữ số mà khi cộng với $8$ ta được kết quả nhỏ hơn $17$.
\choice
{\True $8$}
{$9$}
{$7$}
{$6$}
\loigiai{
Ta có: $x + 8 < 17 \Rightarrow x < 9$. Số tự nhiên lớn nhất có 1 chữ số nhỏ hơn 9 là $8$.
\textbf{Đáp án của câu hỏi là A}
}
\end{ex}

\begin{ex} %[2D2V14]
Trong một hộp có $18$ que tính gồm ba màu xanh, đỏ, vàng. Biết có $7$ que màu xanh và $6$ que màu đỏ. Hỏi có bao nhiêu que tính màu vàng?
\choice
{$6$ que tính}
{$4$ que tính}
{\True $5$ que tính}
{$7$ que tính}
\loigiai{
Số que tính màu vàng: $18 - (7 + 6) = 18 - 13 = 5$ (que tính).
\textbf{Đáp án của câu hỏi là C}
}
\end{ex}

\begin{ex} %[2D2V14]
Hiệu giữa số chẵn lớn nhất có một chữ số và số lẻ bé nhất có một chữ số khác 1 là
\choice
{$6$}
{\True $5$}
{$4$}
{$7$}
\loigiai{
Số chẵn lớn nhất có 1 chữ số là 8. Số lẻ bé nhất khác 1 là 3.\nHiệu: $8 - 3 = 5$.
\textbf{Đáp án của câu hỏi là B}
}
\end{ex}

\begin{ex} %[2D2V14]
Hình vẽ gồm hai đoạn thẳng $AB = 9\text{ cm}$ và $BC = 8\text{ cm}$ ghép lại thành đường gấp khúc $ABC$. Sau đó người ta cắt bớt đi $6\text{ cm}$. Độ dài đường gấp khúc còn lại là
\choice
{$10\text{ cm}$}
{$12\text{ cm}$}
{$9\text{ cm}$}
{\True $11\text{ cm}$}
\loigiai{
Độ dài đường gấp khúc: $9 + 8 = 17\text{ (cm)}$.\nCòn lại: $17 - 6 = 11\text{ (cm)}$.
\textbf{Đáp án của câu hỏi là D}
}
\end{ex}

\begin{ex} %[2D2V14]
Một rổ mận có $16$ quả. Chị ăn $4$ quả, em ăn $5$ quả. Hỏi hai chị em đã ăn hết bao nhiêu quả mận?
\choice
{\True $9$ quả}
{$8$ quả}
{$7$ quả}
{$10$ quả}
\loigiai{
Hai chị em đã ăn hết: $4 + 5 = 9$ (quả mận).
\textbf{Đáp án của câu hỏi là A}
}
\end{ex}

\begin{ex} %[2D2V14]
Tìm hai số có tổng bằng $16$ và hiệu bằng $4$.
\choice
{$9\text{ và } 7$}
{$11\text{ và } 5$}
{\True $10\text{ và } 6$}
{$8\text{ và } 8$}
\loigiai{
Số lớn: $(16 + 4) : 2 = 10$. Số bé: $10 - 4 = 6$.
\textbf{Đáp án của câu hỏi là C}
}
\end{ex}

\textbf{Phần 2. Trắc nghiệm đúng sai}
\setcounterpageref{ex}{0}

\begin{ex} %[2D2V14]
Xét tính đúng sai của các khẳng định ôn tập Chủ đề 2:
\choiceTF
{\True Biểu thức $9 + 7 - 8$ có kết quả bằng 8.}
{Trong phép trừ có nhớ phạm vi 20, số bị trừ luôn nhỏ hơn 10.}
{\True Tổng của số lẻ và số lẻ luôn là số chẵn.}
{Nếu $x + 8 = 15$ thì $x = 8$.}
\loigiai{
\begin{itemize}
    \item Ý a) Đúng: $9 + 7 = 16$; $16 - 8 = 8$.
    \item Ý b) Sai: Số bị trừ luôn lớn hơn hoặc bằng 11.
    \item Ý c) Đúng: Ví dụ: $7 + 9 = 16$ (số chẵn).
    \item Ý d) Sai: $x = 15 - 8 = 7$ (không phải 8).
\end{itemize}
\begin{itemize}
    \item a) \textbf{Đúng}.
    \item b) \textbf{Sai}.
    \item c) \textbf{Đúng}.
    \item d) \textbf{Sai}.
\end{itemize}
}
\end{ex}

\begin{ex} %[2D2V14]
Một bác nông dân nuôi 8 con thỏ trắng và 9 con thỏ nâu:
\choiceTF
{\True Bác nuôi tất cả 17 con thỏ.}
{\True Số thỏ nâu nhiều hơn số thỏ trắng 1 con.}
{Mỗi con thỏ có 2 cái tai và 2 cái chân.}
{\True Tổng số tai của cả đàn thỏ là 34 cái tai.}
\loigiai{
\begin{itemize}
    \item Ý a) Đúng: $8 + 9 = 17$ con thỏ.
    \item Ý b) Đúng: $9 - 8 = 1$ con.
    \item Ý c) Sai: Mỗi con thỏ có 2 cái tai và 4 cái chân.
    \item Ý d) Đúng: $17 \times 2 = 34$ cái tai. Đúng.
\end{itemize}
\begin{itemize}
    \item a) \textbf{Đúng}.
    \item b) \textbf{Đúng}.
    \item c) \textbf{Sai}.
    \item d) \textbf{Đúng}.
\end{itemize}
}
\end{ex}

\textbf{Phần 3. Trả lời ngắn}
\setcounterpageref{ex}{0}

\begin{ex} %[2D2V14]
Tính giá trị của biểu thức: $16 - 8 + 7$.
\shortans{15}
\loigiai{
\textbf{Đáp số:} $15$\\
$16 - 8 = 8$; $8 + 7 = 15$.
}
\end{ex}

\begin{ex} %[2D2V14]
Tìm một số biết rằng nếu lấy số đó cộng với 9 rồi trừ đi 6 thì được 11.
\shortans{8}
\loigiai{
\textbf{Đáp số:} $8$\\
Số đó là: $11 + 6 - 9 = 8$.
}
\end{ex}

\begin{ex} %[2D2V14]
Có 17 quả bóng bay, bị nổ mất 8 quả, sau đó được thổi thêm 5 quả. Hỏi có tất cả bao nhiêu quả bóng bay?
\shortans{14}
\loigiai{
\textbf{Đáp số:} $14$\\
$17 - 8 + 5 = 14$ quả.
}
\end{ex}

\begin{ex} %[2D2V14]
Hai bạn có tổng cộng 18 cái kẹo. Bạn thứ nhất hơn bạn thứ hai 4 cái kẹo. Hỏi bạn thứ nhất có bao nhiêu cái kẹo?
\shortans{11}
\loigiai{
\textbf{Đáp số:} $11$\\
Số kẹo của bạn thứ nhất: $(18 + 4) : 2 = 11$ cái.
}
\end{ex}

\begin{ex} %[2D2V14]
Tìm số tự nhiên x biết: $13 - x = 11 - 6$.
\shortans{8}
\loigiai{
\textbf{Đáp số:} $8$\\
$11 - 6 = 5 \Rightarrow x = 13 - 5 = 8$.
}
\end{ex}

\begin{ex} %[2D2V14]
Tìm tổng của số liền sau số 8 và số liền trước số 10.
\shortans{18}
\loigiai{
\textbf{Đáp số:} $18$\\
Số liền sau số 8 là 9. Số liền trước số 10 là 9. Tổng: $9 + 9 = 18$.
}
\end{ex}
